\documentclass{article}
\usepackage{hyperref}
\usepackage{Style}

\nocite{*} % Comentar si quiero citar
%\addbibresource{bibliografia.bib} % Quitar el comentado si quiero usar bibliografia

%% ----- Comando especial -----
\def \P {\mathbb{P}}

\begin{document}

\begin{minipage}{2.5cm}
    \includegraphics[width=2cm]{imagen_puc.jpg}
\end{minipage}
\begin{minipage}{12cm}
    {\sc Pontificia Universidad Católica de Chile\\
    Facultad de Matemáticas\\
    Departamento de Matemática\\
    Professor: Jesse Huang -- Student: Benjamín Mateluna}
\end{minipage}
\vspace{2ex}

{\centerline{\bf Selected Topics on Complex Algebraic Varieties - MPG3320 III}
\centerline{\bf Assignment 2}}
\centerline{\bf September 28, 2026}

\vspace{8mm}

\vspace{2mm}
\noindent\textbf{Lemma 1:} \emph{Let $V\subseteq\A_{k}^{n}$ be an algebraic variety. Then there
exist irreducible affine varieties $V_{1},\dots,V_{m}$ such that}
\begin{equation*}
    V=\bigcup V_{i} \hhtext{y} V_{i}\not\subset V_{j} \htext{for all}i\neq j
\end{equation*}

\begin{proof}
    We define
    \begin{equation*}
        \mathcal{C}:=\{
            V\subseteq\A_{k}^{n}\text{ affine variety}:V\text{ is not an union of irreducibles}
        \}
    \end{equation*}
    If $\mathcal{C}$ is empty there is nothing to prove, otherwise, there exists a minimal 
    $V\in\mathcal{C}$, and by definition of $\mathcal{C}$, it is reducible. This contradicts its
    minimality.

    Let $V=\bigcup V_{i}$ with $V_{i}$ irreducible and suppose that $V_{i}\not\subset V_{j}$ for 
    all $i\neq j$. Suppose that
    \begin{equation*}
        \bigcup V_{i}=\bigcup W_{j} 
        \htext{where}W_{i}\not\subset W_{j} \hhtext{and} V_{i},W_{j}\neq\emptyset
    \end{equation*}
    Notice that $V_{1}=V_{1}\cap V=\bigcup(V_{1}\cap W_{j})$, since $V_{1}$ is irreducible, there
    exists a unique $j$ such that $V_{1}=V_{1}\cap W_{j}$, that is, $V_{1}\subseteq W_{j}$. In the
    same way, there exists a unique $i$ such that $W_{j}\subseteq V_{i}$, which implies that
    $V_{1}\subseteq V_{i}$, then $i=1$ and thus $V_{1}=W_{j}$.
\end{proof}

Since $k[x_{1},\dots,x_{n}]$ is noetherian, every collection of algebraic varieties has a minimal
element.

\vspace{2mm}
\noindent\textbf{Lemma 2:} \emph{Let $X\subseteq\A_{k}^{n}$ be an algebraic variety and 
$X=\bigcup X_{i}$ its decomposition into irreducibles, then}
\begin{equation*}
    dim(X)=max \text{ } dim(X_{i})
\end{equation*}

\begin{proof}
    Given $V\subseteq X$ an irreducible affine variety, we see that $V=\bigcup(V\cap X_{i})$ and
    since $V$ is irreducible it follows that $V\subseteq X_{j}$ for some $j$.

    \vspace{2mm}
    Let $V_{0}\subset V_{1}\subset\dots\subset V_{r}$ be irreducible varieties in $X$, by the 
    previous observation, we have that $V_{i}\subseteq X_{j}$ for some $j$ and for all $i$. Then,
    \begin{equation*}
        r\leq dim(X_{j})\leq max \text{ } dim(X_{i})
        \hhtext{hence}
        dim(X)\leq max \text{ } dim(X_{i})
    \end{equation*}
    On the other hand, if $V_{0}\subset V_{1}\subset\dots\subset V_{r}$ are irreducible varieties
    in $X_{i}$, in particular, they are so in $X$, then
    \begin{equation*}
        r\leq dim(X)
        \hhtext{and therefore}
        dim(X_{i})\leq dim(X)
    \end{equation*}
    which implies that $max \text{ } dim(X_{i})\leq dim(X)$.
\end{proof}

Both results hold for projective varieties, using identical arguments.

\vspace{3mm}
\noindent\textbf{Problem 1.} Let us look the following lemma

\vspace{2mm}
\noindent\textbf{Lemma 1.1:} \emph{Let $f:X\to Y$ be a surjective morphism between irreducible 
algebraic varieties, then $dim(Y)\leq dim(X)$}.

\begin{proof}
    Let us see that the $k-$algebra morphism $f^{*}:k[Y]\to k[X]$ is injective. Let 
    $g\in\kr{f^{*}}$, then $(g\circ f)(x)=0$ for all $x\in X$, in other words,
    \begin{equation*}
        Y=f(X)\subseteq V(g) \hhtext{then} g|_{Y}\equiv0
    \end{equation*}
    Therefore $k(Y)\subseteq k(X)$ and thus $dim(Y)=trdeg_{k}(k(Y))\leq trdeg_{k}(k(X))=dim(X)$.
\end{proof}

\newpage
Let $X=\bigcup X_{i}$ and $Y=\bigcup Y_{j}$ be the decomposition into irreducibles of $X$ and $Y$
respectively. Notice that
\begin{equation*}
    X\times Y=\bigcup X_{i}\times\bigcup Y_{j}=\bigcup(X_{i}\times Y_{j})
\end{equation*}
and since the product of irreducible varieties is irreducible, it follows that 
$\bigcup(X_{i}\times Y_{j})$ is the decomposition into irreducibles of $X\times Y$. Let $X_{i}$ 
and $Y_{j}$ of maximal dimension, by Noether's normalization lemma there exist surjective 
morphisms
\begin{equation*}
    f:X_{i}\to\A_{k}^{\dim(X_{i})} \hhtext{and} g:Y_{j}\to\A_{k}^{\dim(Y_{j})}
\end{equation*}
We have a surjective morphism between irreducible varieties
$\varphi=(f,g):X_{i}\times Y_{j}\to\A_{k}^{\dim(X_{i})+\dim(Y_{j})}$, then
\begin{equation*}
    dim(X)+dim(Y)=dim(X_{i})+dim(Y_{j})\leq dim(X\times Y)
\end{equation*}
On the other hand, let $X_{i}\times X_{j}$ be of maximal dimension. By by Noether's normalization 
lemma, $k[X_{i}]$ is algebraically generated by $dim(X_{i})$ elements, and similarly, $k[Y_{j}]$ 
is algebraically generated by $dim(Y_{j})$ elements.

\vspace{2mm}
Therefore, $k[X_{i}\times Y_{j}]\cong k[X_{i}]\otimes_{k} k[Y_{j}]$ is algebraically generated by
$dim(X_{i})+dim(Y_{j})$ elements. Thus, we see that
\begin{align*}
    dim(X\times Y) &= dim(X_{i}\times Y_{j})=trdeg_{k}(k(X_{i}\times Y_{j})) \\
    &= trdeg_{k}(Frac(k[X_{i}]\otimes_{k} k[Y_{j}]))\leq dim(X_{i})+dim(Y_{j})\leq dim(X)+dim(Y)
\end{align*}
Which finalizes the proof.

\vspace{5mm}
\noindent\textbf{Problem 2.} Let us see $\varphi$ is surjective. Let $[a]\in \P_{k}^{n}$, we 
consider the hyperplane $V^{p}(h_{a})\in(\P_{k}^{n})^{\lor}$, where 
$h_{a}=a_{0}x_{0}+\dots+a_{n}x_{n}$, note that $\varphi(V^{p}(h_{a}))=[a]$. Let $[a]=[b]$, then
there exists $\lambda\neq0$ such that $a=\lambda b$, hence
\begin{equation*}
    h_{a}=\sum_{i=0}^{n}a_{i}x_{i}=\sum_{i=0}^{n}\lambda b_{i}x_{i}
    =\lambda\sum_{i=0}^{n}b_{i}x_{i}=\lambda h_{b}
\end{equation*}
so $h_{a}$ and $h_{b}$ differ by a non-zero scalar, therefore they define the same hyperplane, in 
other words $V^{p}(h_{a})=V^{p}(h_{b})$. We define
\begin{equation*}
    X=\left\{V^{p}(h_{a}):[1:0:\dots:0]\in V^{p}(h_{a})\right\}
\end{equation*}
and $V=\varphi(X)$. We claim that $V=V^{p}(a_{0})$. Indeed, if $[a]\in V^{p}(a_{0})$ then 
$a_{0}=0$ and therefore
\begin{equation*}
    h_{a}=\sum_{i=1}^{n}a_{i}x_{i} \hhtext{then} h_{a}(1,0,\dots,0)=0
\end{equation*}
and thus $[a]\in V$. On the other hand, if $V^{p}(h_{a})\in X$ we have that 
$0=h_{a}(1,0,\dots,0)=a_{0}$, then $[a]\in V^{p}(a_{0})$. Which concludes what was requested.

\vspace{5mm}
\noindent\textbf{Problem 3.}
\begin{enumerate}
    \item[(a)] Since the ideal $I=(y^{2}-x^{3}-ax-b)$ is generated by a single element, its 
    homogenization is equal to the ideal $(y^{2}z-x^{3}-axz^{2}-bz^{3})$, then
    \begin{equation*}
        \overline{V(I)}=V^{p}(\widetilde{I})=V^{p}(y^{2}z-x^{3}-axz^{2}-bz^{3})
    \end{equation*}
    Then, if $z=0$ we have that $x=0$ and therefore $y=1$. We conclude that the projective closure
    of $V(I)$ added a unique point at infinity, which corresponds to $[0:0:1]$.

    \item [(b)] We want to show that $\{x_{1}^{2}-x_{2}^{2}-1,x_{3}-x_{1}\}$ is a Gröbner basis 
    for $I=(x_{1}^{2}-x_{2}^{2}-1,x_{3}-x_{1})$ with the lexicographic monomial order given by
    $x_{3}\geq x_{1}\geq x_{2}$. Under this order,
    \begin{equation*}
        LM_{\leq}(x_{1}^{2}-x_{2}^{2}-1)=x_{1}^{2} \hhtext{y}
        LM_{\leq}(x_{3}-x_{1})=x_{3}
    \end{equation*}
    It suffices to show that, given $f\in I$, its leading monomial is divisible by $x_{3}$ or
    $x_{2}^{2}$. Since $f\in I$, there exist $g,h\in k[x_{1},x_{2},x_{3}]$ such that
    \begin{equation*}
        f=g\sbullet(x_{1}^{2}-x_{2}^{2}-1)+h\sbullet(x_{3}-x_{1})
    \end{equation*}
    We have two cases: if the leading monomial of $f$ has at least one factor $x_{3}$, then it is
    divisible by $x_{3}$. If it does not, then no monomial contains it as a factor; otherwise it
    would be the leading monomial. Thus, $x_{3}$ is a free variable, so
    \begin{equation*}
        f(x_{1},x_{2},x_{3})
        =f(x_{1},x_{2},x_{1})=g(x_{1},x_{2},x_{1})\sbullet(x_{1}^{2}-x_{2}^{2}-1)
    \end{equation*}
    Then, the leading monomial of $f$ is the leading monomial of $g$ multiplied by $x_{1}^{2}$, 
    since the order respects multiplication. Therefore, 
    $\overline{W}=V^{p}(x_{1}^{2}-x_{2}^{2}-x_{0}^{2},x_{3}-x_{1})$. If $x_{0}=0$, we have
    \begin{equation*}
        x_{1}^{2}=x_{2}^{2} \hhtext{then} x_{1}=\pm x_{2}
    \end{equation*}
    Since $x_{3}=x_{1}$, it follows that $x_{1}\neq0$. Therefore, the points at infinity are
    \begin{equation*}
        [0:1:1:1] \hhtext{and} [0:1:-1:1]
    \end{equation*}
\end{enumerate}

\vspace{5mm}
\noindent\textbf{Problem 4.} The twisted cubic is given by the ideal $(y-x^{2},z-x^{3})$, that is,
$Y=V(y-x^{2},z-x^{3})$ and it is parameterized by the expression $(t,t^{2},t^{3})$. We see that
\begin{equation*}
    k[x,y,z]/(y-x^{2},z-x^{3})\cong k[x]
\end{equation*}
and therefore $(y-x^{2},z-x^{3})$ is a prime ideal, so by the Nullstellensatz we have that
\begin{equation*}
    I(Y)=\sqrt{(y-x^{2},z-x^{3})}=(y-x^{2},z-x^{3})=:I
\end{equation*}
Let $J=(yw-x^{2},xz-y^{2},zw-xy)$ and note that $J\subseteq\widetilde{I}$, since the generators of
$J$ are in $\widetilde{I}$. It follows that $\overline{Y}=V^{p}(\widetilde{I})\subseteq V^{p}(J)$,
we would like to prove that $V^{p}(J)=\overline{Y}$. Consider the morphism
\begin{align*}
    \phi:\P_{k}^{1} &\to V^{p}(J) \\
    [s:t] &\mapsto [s^{3}:s^{2}t:st^{2}:t^{3}]
\end{align*}
Evaluating at the generators of $J$ we see that $\phi$ is well-defined, moreover, its inverse is 
given by
\begin{equation*}
    [w:x:y:z]\mapsto\begin{cases}
        [w:x] \quad \text{si $w\neq0$ ó $x\neq0$} \\
        [y:z] \quad \text{si $y\neq0$ ó $z\neq0$}
    \end{cases}
\end{equation*}
Therefore, $\phi$ is a homeomorphism, in particular, it is a closed map. Since 
$\phi(U_{0})\subseteq Y$, we have the following
\begin{equation*}
    V^{p}(J)=\phi(\P_{k}^{1})=\phi(\overline{U_{0}})=\overline{\phi(U_{0})}\subseteq\overline{Y}
\end{equation*}
Suppose, by contradiction, that $\widetilde{I}=(yw-x^{2},zw^{2}-x^{3})$, then 
$V^{p}(yw-x^{2},zw^{2}-x^{3})=\overline{Y}=V^{p}(J)$. The parameterization $\phi$ tell us that 
there is a unique point $Y$ such that $w=0$, namely, $[0:0:0:1]$. However, in the variety
$V^{p}(yw-x^{2},zw^{2}-x^{3})$ we have a copy of $\P_{k}^{1}$, since $x=0$ and $y,z$ remain free.
This contradiction proves the claim.

\vspace{5mm}
\noindent\textbf{Problem 5.} No answer.

\vspace{5mm}
\noindent\textbf{Problem 6.}
\begin{enumerate}
    \item[(a)] Clearly the projection is a morphism of algebraic varieties, let us see that it is
    finite and surjective. In the first place, note that $k[\A_{k}^{2}]=k[x,y]$, then the induced
    morphism
    \begin{equation*}
        \pi^{*}:k[x,y]\to k[X]
    \end{equation*}
    is such that $x\mapsto x$ and $y\mapsto y$. Furthermore, given $p\in k[x,y,z]$ we have that
    $deg_{z}(p(z^{3}+xz-y))\geq1$, which implies that $\pi^{*}$ is injective. Moreover, the 
    element $z$ satisfies the equation
    \begin{equation*}
        z^{3}+xz-y=0
    \end{equation*}
    so $k[X]$ is generated by $\{1,z,z^{2}\}$ as a $k[x,y]-$module, that is, it is finitely 
    generated. Therefore, $k[X]$ is an integral entension of $k[x,y]$. This proves that $\phi$ is
    a finite, surjective morphism.

    \item[(c)] Since $k$ is a field of characteristic $0$, we have $\Q\subseteq k$. Let 
    $(a,b)\in\A_{k}^{2}$, then
    \begin{equation*}
        \pi^{-1}(a,b)=\{(a,b,z)\in\A_{k}^{3}:z^{3}+az-b\}
    \end{equation*}
    Then $z$ corresponds to the root of a cubic equation, whose solutions are given by the 
    Cardano-Tartaglia formula, namely
    \begin{equation*}
        z_{1}=\sqrt[3]{\frac{b}{2}+\sqrt{\frac{b^{2}}{4}+\frac{a^{3}}{27}}}
        +\sqrt[3]{\frac{b}{2}-\sqrt{\frac{b^{2}}{4}+\frac{a^{3}}{27}}}=u+v
    \end{equation*}
    The other two solutions are $z_{2}=\omega u+\omega^{2} v$ and $z_{3}=\omega^{2} u+\omega v$,
    where $\omega$ is a cube root of unity distinct from $1$. If $(a,b)$ are such that
    \begin{equation*}
        \frac{b^{2}}{4}+\frac{a^{3}}{27}\neq0
    \end{equation*}
    so $\pi^{-1}(a,b)$ consists of three distinct elements. Indeed, if $z_{1}=z_{2}$ then we have
    \begin{equation*}
        0=z_{1}-z_{2}=(1-\omega)u+(1-\omega^{2})v
    \end{equation*}
    which implies that
    \begin{align*}
        0 &= (1-\omega)^{3}u^{3}+(1-\omega^{2})^{3}v
        =-3\omega(1-\omega)u^{3}-3\omega^{2}(1-\omega^{2})v^{3} \\
        &= -3\omega(1-\omega)u^{3}+3\omega(1-\omega)v^{3}=3\omega(1-\omega)(v^{3}-u^{3})
    \end{align*}
    Therefore, $u^{3}=v^{3}$, it follows that $\frac{b^{2}}{4}+\frac{a^{3}}{27}=0$, this prove 
    that $z_{1}\neq z_{2}$. The other cases are obtained similarly. Thus, $\pi$ is a covering map
    of degree $3$ outside of
    \begin{equation*}
        B=V\left(\frac{b^{2}}{4}+\frac{a^{3}}{27}\right)
    \end{equation*}
    Then, the branch locus of $\pi$ corresponds to $B$.

    \item[(b)] Since $\Q\subseteq k$, we can write
    \begin{equation*}
        \omega=-\frac{1}{2}+\frac{\sqrt{3}}{2}i \hhtext{and} 
        \omega^{2}=-\frac{1}{2}-\frac{\sqrt{3}}{2}i \htext{where $i$ is a root of -1}
    \end{equation*}
    In $B$ we have two options, $(a,b)=0$ ó $(a,b)\neq0$ in which case $a\neq0$ and $b\neq0$. In
    the first one we have
    \begin{equation*}
        \pi^{-1}(0,0)=(0,0,0)
    \end{equation*}
    and in the second case, replacing the value of $\omega$ and $\omega^{2}$ into $z_{2}$ and 
    $z_{3}$, we obtain the following
    \begin{equation*}
        \pi^{-1}(a,b)=\left\{
            \left(a,b,2\sqrt[3]{b/2}\right),\left(a,b,-\sqrt[3]{b/2}\right)
        \right\}
    \end{equation*}
    Thus,
    \begin{equation*}
        R=\pi^{-1}(B)=\{0\}\cup\left\{
            \left(a,b,2\sqrt[3]{b/2}\right),\left(a,b,-\sqrt[3]{b/2}\right):(a,b)\in B
        \right\}
    \end{equation*}

    \item[(d)] Outside of $B$, the fiber of a point has $3$ distinct elements and each with 
    multiplicity $1$. In $B$ we have two cases, the origin, in which case the fiber consists of a 
    unique point with multiplicity $3$. And in the second case, we have two distinct points in
    the fiber, one with multiplicity $1$ $(z_{1})$ and another with multiplicity $2$ 
    $(z_{2}=z_{3})$.
\end{enumerate}

\vspace{5mm}
\noindent\textbf{Problem 7.} The segre embedding of $\P_{k}^{1}\times\P_{k}^{1}$ into $\P_{k}^{3}$
is given by the map
\begin{align*}
    \sigma_{1,1}:\P_{k}^{1}\times\P_{k}^{1} &\to \P_{k}^{3} \\
    ([x:y],[a:b]) &\mapsto [xa:xb:ya:yb]
\end{align*}
Whose image corresponds to the projective variety given by the equation $XW-YZ$. Note that the
expression is equivalent to the zero locus of a quadratic form, indeed,
\begin{equation*}
    XW-YZ=[X:Y:Z:W]\begin{pmatrix}
        0 & 0 & 0 & 1/2 \\
        0 & 0 & -1/2 & 0 \\
        0 & -1/2 & 0 & 0 \\
        1/2 & 0 &  0 & 0
    \end{pmatrix}[X:Y:Z:W]^{T}
\end{equation*}
Let $B$ be the matrix associated with such quadratic form, we see that $rank(B)=4$. As the matrix 
$A$, which defines the quadratic form $Q$, has rank 4, by linear algebra there exists $P$ 
invertible such that $P^{T}AP=B$. We define
\begin{align*}
    \varphi:\P_{k}^{3} &\to \P_{k}^{3} \\
    [x] &\mapsto [P^{-1}x]
\end{align*}
Let us see that $\varphi(X)=\Sigma_{1,1}$. Let $[x]\in X$, then
\begin{equation*}
    \varphi(x)^{T}B\varphi(x)=(P^{-1}x)^{T}P^{T}AP(P^{-1}x)=x^{T}(P^{-1})^{T}P^{T}AP(P^{-1}x)
    =x^{T}Ax=0
\end{equation*}
and therefore $\varphi([x])\in\Sigma_{1,1}$. On the other hand, let $[y]\in\Sigma_{1,1}$ and we 
consider $[x]=[Py]$, we have
\begin{equation*}
    x^{T}Ax=(Py)^{T}APy=y^{T}P^{T}APy=y^{T}By=0
\end{equation*}
that is, $[y]\in\varphi(X)$. Which concludes the proof.

%\printbibliography % Quitar el comentado si quiero usar bibliografia

\end{document}
